\documentclass[12pt,a4paper]{article}
\usepackage[utf8]{inputenc}
\usepackage[T1]{fontenc}
\usepackage[brazil]{babel}
\usepackage{amsmath,amssymb}
\usepackage{booktabs}
\usepackage{graphicx}
\usepackage{geometry}
\usepackage{caption}
\usepackage{placeins}
\geometry{margin=2.5cm}

\graphicspath{{trabalho3/figures/}}

\title{Segundo Trabalho de Reconhecimento de Padr\~oes}
\author{Jakson da Rocha Almeida}
\date{12 de junho de 2026}

\begin{document}

\maketitle

\section*{Introdu\c{c}\~ao}

Este relat\'orio documenta a implementa\c{c}\~ao das quest\~oes do segundo trabalho, realizada em Python com \texttt{numpy}, \texttt{matplotlib} e \texttt{scikit-learn}. Os dados seguem duas classes Gaussianas bidimensionais com a mesma matriz de covari\^ancia
$\boldsymbol{\Sigma}=\bigl[\begin{smallmatrix}1&0{,}6\\0{,}6&1\end{smallmatrix}\bigr]$,
m\'edias $\boldsymbol{\mu}_1=(0,0)^\top$ e $\boldsymbol{\mu}_2=(2,0)^\top$, e probabilidades \`a priori iguais.

Para tornar os sorteios reproduz\'{\i}veis (sem contrariar o enunciado, que pede inteiros em $[0,500]$ como ra\'{\i}zes), utilizamos sementes mestras fixas no notebook:
\texttt{MASTER\_SEED}=2027 para obter a semente do treino da Quest\~ao~2 ($316$)
e \texttt{MASTER\_Q34}=2028 para a ``raiz'' das Quest\~oes~3--4, correspondendo ao inteiro $283$ usado na gera\c{c}\~ao do teste.
As figuras s\~ao gravadas na pasta \texttt{documents/\allowbreak trabalho2/\allowbreak figures/}.

\section*{Quest\~ao 1 --- Defini\c{c}\~ao das classes}

Duas classes $C_1$ e $C_2$, cada uma $\mathcal{N}(\boldsymbol{\mu}_i,\boldsymbol{\Sigma})$ com $\boldsymbol{\Sigma}$ compartilhada, conforme o enunciado.

\section*{Quest\~ao 2 --- Treino e fronteira de Bayes}

Sorteou-se o inteiro $316 \in [0,500]$ e geraram-se $n=1000$ amostras por classe.
Superp\~oe-se a \textbf{fronteira \'otima de Bayes} (linear neste modelo), igual \`a derivada no primeiro trabalho para esta mesma configura\c{c}\~ao:
\[
h(\mathbf{x}) = 3{,}125\,x_1 - 1{,}875\,x_2 - 3{,}125 = 0.
\]

\begin{figure}[ht]
\centering
\includegraphics[width=0.62\textwidth]{questao2_treino_fronteira_bayes_t1.png}
\caption{Quest\~ao~2: treino ($1000$ pontos por classe) e fronteira de Bayes.}
\label{fig:q2}
\end{figure}

\section*{Quest\~oes 3 e 4 --- Conjunto de teste}

O inteiro sorteado em $[0,500]$ para a ra\'{\i}z do teste foi $283$.
Geraram-se novamente $1000$ amostras por classe, independentes do treino, totalizando $2000$ pontos de teste etiquetados ($y\in\{0,1\}$ com $C_1\to 0$, $C_2\to 1$).

\section*{Erro te\'orico de Bayes (refer\^encia para as compara\c{c}\~oes)}

Com covari\^ancia comum e priors iguais, a regra de Bayes \'e linear e a probabilidade de erro vale
$P_e^{\mathrm{(Bayes)}}=\Phi(-d_M/2)$,
onde $d_M=\sqrt{(\boldsymbol{\mu}_2-\boldsymbol{\mu}_1)^\top\boldsymbol{\Sigma}^{-1}(\boldsymbol{\mu}_2-\boldsymbol{\mu}_1)}$.
Aqui $d_M=2{,}5$, logo $P_e^{\mathrm{(Bayes)}}\approx 0{,}1056$
aproximadamente $10{,}6\%$.

\section*{Quest\~ao 5 --- Perceptron (treino completo)}

Treinou-se um perceptron linear (\texttt{sklearn.linear\_model.Perceptron}) com todo o conjunto de treino, crit\'erios de parada via \texttt{max\_iter} e \texttt{tol}, e avaliou-se no teste.

\paragraph{Resultado.} Erro emp\'{\i}rico estimado $\hat P_e \approx 0{,}1655$.
Esperava-se valor \textbf{superior} ao de Bayes: o perceptron minimiza erros por atualiza\c{c}\~oes e n\~ao incorpora explicitamente $\boldsymbol{\Sigma}^{-1}$ nem maximiza verossimilhan\c{c}a das Gaussianas.

\begin{figure}[ht]
\centering
\includegraphics[width=0.62\textwidth]{questao5_perceptron_teste.png}
\caption{Quest\~ao~5: perceptron no conjunto de teste e refer\^encia Bayes (verde).}
\label{fig:q5}
\end{figure}

\section*{Quest\~oes 6 e 7 --- Perceptron com $50{+}50$ amostras}

Do treino completo, extra\'{\i}ram-se dois blocos disjuntos de $50$ pontos por classe (permuta\c{c}\~oes fixas com \texttt{rng\_sub = np.random.default\_rng(3030)}).

\paragraph{Resultados.}
\begin{itemize}
  \item Quest\~ao~6 (primeiro bloco): $\hat P_e \approx 0{,}3565$.
  \item Quest\~ao~7 (segundo bloco): $\hat P_e \approx 0{,}1080$.
\end{itemize}
A dispers\~ao entre os dois blocos \'e grande: com poucos exemplos e overlap das classes, a reta aprendida depende fortemente do subconjunto.

\begin{figure}[ht]
\centering
\includegraphics[width=0.62\textwidth]{questao6_perceptron_50_teste.png}
\caption{Quest\~ao~6: perceptron com $50{+}50$ (bloco~1).}
\label{fig:q6}
\end{figure}

\begin{figure}[ht]
\centering
\includegraphics[width=0.62\textwidth]{questao7_perceptron_50_teste.png}
\caption{Quest\~ao~7: perceptron com $50{+}50$ (bloco~2).}
\label{fig:q7}
\end{figure}

\section*{Quest\~oes 8 a 10 --- M\'{\i}nimos quadrados (pseudoinversa)}

Implementou-se o classificador linear por MQ com matriz aumentada $\tilde{\boldsymbol{X}}=[\mathbf{1}\mid \boldsymbol{X}]$, alvos $\pm 1$ e solu\c{c}\~ao $\hat{\boldsymbol{\theta}}=\tilde{\boldsymbol{X}}^+\boldsymbol{y}$ (\texttt{numpy.linalg.pinv}), classe prevista pela polaridade de $\tilde{\boldsymbol{x}}^\top\hat{\boldsymbol{\theta}}$.

\paragraph{Resultados.}
\begin{itemize}
  \item Quest\~ao~8 (treino completo): $\hat{\boldsymbol{\theta}}\approx[-0{,}599,\;0{,}612,\;-0{,}396]^\top$ (formato $[b,w_1,w_2]$), $\hat P_e \approx 0{,}1035$.
  \item Quest\~ao~9 (mesmo $50{+}50$ da Q6): $\hat P_e \approx 0{,}1110$.
  \item Quest\~ao~10 (mesmo $50{+}50$ da Q7): $\hat P_e \approx 0{,}1030$.
\end{itemize}

\begin{figure}[ht]
\centering
\includegraphics[width=0.62\textwidth]{questao8_ls_pinv_treino_completo_teste.png}
\caption{Quest\~ao~8: MQ com treino completo.}
\label{fig:q8}
\end{figure}

\begin{figure}[ht]
\centering
\includegraphics[width=0.48\textwidth]{questao9_ls_pinv_50_teste.png}\quad
\includegraphics[width=0.48\textwidth]{questao10_ls_pinv_50_teste.png}
\caption{Quest\~oes~9 e~10: MQ com blocos $50{+}50$.}
\label{fig:q910}
\end{figure}

\section*{Quest\~oes 11 a 13 --- SVM linear com margem suave}

Usou-se \texttt{SVC(kernel="linear",\,C=1.0)}. No treino completo h\'a muitos vetores de suporte porque as classes se sobrep\~oem na distribui\c{c}\~ao geradora.

\paragraph{Resultados.}
\begin{itemize}
  \item Quest\~ao~11: $525$ vetores de suporte no total, $\hat P_e \approx 0{,}1005$.
  \item Quest\~ao~12 (bloco~1): $32$ vetores de suporte, $\hat P_e \approx 0{,}1140$.
  \item Quest\~ao~13 (bloco~2): $33$ vetores de suporte, $\hat P_e \approx 0{,}1050$.
\end{itemize}

\begin{figure}[ht]
\centering
\includegraphics[width=0.62\textwidth]{questao11_svm_linear_teste.png}
\caption{Quest\~ao~11: SVM linear ($C=1$), treino completo.}
\label{fig:q11}
\end{figure}

\begin{figure}[ht]
\centering
\includegraphics[width=0.48\textwidth]{questao12_svm_linear_50_teste.png}\quad
\includegraphics[width=0.48\textwidth]{questao13_svm_linear_50_teste.png}
\caption{Quest\~oes~12 e~13: SVM nos dois blocos $50{+}50$.}
\label{fig:q1213}
\end{figure}

\FloatBarrier
\section*{Quadro-resumo --- probabilidade de erro no teste}

\begin{table}[ht]
\centering
\begin{tabular}{@{}llcc@{}}
\toprule
\textbf{Quest\~ao} & \textbf{M\'etodo} & \textbf{Tamanho do treino} & $\hat P_e$ \\
\midrule
--- & Bayes (te\'orico) & --- & $\approx 0{,}1056$ \\
5 & Perceptron & $1000{+}1000$ & $0{,}1655$ \\
6 & Perceptron & $50{+}50$ (bloco~1) & $0{,}3565$ \\
7 & Perceptron & $50{+}50$ (bloco~2) & $0{,}1080$ \\
8 & MQ (pseudoinversa) & $1000{+}1000$ & $0{,}1035$ \\
9 & MQ & $50{+}50$ (bloco~1) & $0{,}1110$ \\
10 & MQ & $50{+}50$ (bloco~2) & $0{,}1030$ \\
11 & SVM linear ($C{=}1$) & $1000{+}1000$ & $0{,}1005$ \\
12 & SVM linear & $50{+}50$ (bloco~1) & $0{,}1140$ \\
13 & SVM linear & $50{+}50$ (bloco~2) & $0{,}1050$ \\
\bottomrule
\end{tabular}
\caption{Erros emp\'{\i}ricos no mesmo conjunto de teste ($2000$ amostras).}
\label{tab:resumo}
\end{table}

\section*{Quest\~ao 14 --- S\'{\i}ntese e compara\c{c}\~ao com Bayes}

O limiar \'otimo de Bayes ($\approx 10{,}6\%$ de erro) serve como refer\^encia inferior para este problema gerador.

Com treino completo, tanto \textbf{m\'{\i}nimos quadrados} ($10{,}35\%$) quanto o \textbf{SVM linear} com $C{=}1$ ($10{,}05\%$) aproximaram-se desse valor, o que \'e esperado: a fronteira de Bayes \'e linear e ambos os m\'etodos produzem separadores lineares adequados a dados quase linearmente separ\'aveis em m\'edia.

O \textbf{perceptron} no treino completo ficou substancialmente acima ($16{,}55\%$), consistente com o facto de o algoritmo n\~ao estimar directamente o classificador likelihood-ratio \'otimo para Gaussianas.

Com apenas $50{+}50$ exemplos, todos os m\'etodos oscilam em torno do valor de Bayes, mas o \textbf{perceptron} mostrou o comportamento mais extremo (de $\approx 35{,}7\%$ a $\approx 10{,}8\%$ entre os dois blocos), evidenciando forte sensibilidade ao ru\'{\i}do amostral.
\textbf{MQ} e \textbf{SVM} mantiveram diferen\c{c}as menores entre blocos ($\approx 0{,}8$ p.p.\ e $\approx 0{,}9$ p.p., respectivamente), mais coerentes com modelos que estabilizam o hiperplano (solu\c{c}\~ao fechada do MQ; maximiza\c{c}\~ao de margem suavizada no SVM).

\section*{Quest\~ao 15 --- Polariza\c{c}\~ao (vi\'es) e vari\^ancia entre os tr\^es m\'etodos lineares}

Neste trabalho n\~ao estim\'amos vi\'es e vari\^ancia por Monte Carlo com muitas repeti\c{c}\~oes; interpretamos os conceitos de forma \textbf{qualitativa e emp\'{\i}rica} com base em (i)~erro com treino grande e (ii)~dispers\~ao entre os dois subconjuntos fixos de $50{+}50$.

\begin{itemize}
  \item \textbf{Vari\^ancia entre os dois blocos pequenos:} o \textbf{perceptron} apresentou claramente a maior dispers\~ao de $\hat P_e$.
  \textbf{M\'{\i}nimos quadrados} e \textbf{SVM} tiveram dispers\~oes semelhantes e bem menores.
  \item \textbf{Polariza\c{c}\~ao relativa a Bayes com treino completo:} o perceptron ficou mais afastado do erro te\'orico de Bayes.
  O SVM obteve $\hat P_e$ mais pr\'oximo de $P_e^{\mathrm{(Bayes)}}$, seguido de perto pelo MQ.
\end{itemize}

Assim, \textbf{n\~ao h\'a um \'unico} dos tr\^es que minimize simultaneamente, nos dois crit\'erios observados, polariza\c{c}\~ao e vari\^ancia: se se privilegia estabilidade entre os dois blocos $50{+}50$, destacam-se \textbf{MQ} e \textbf{SVM}; se se privilegia proximidade ao Bayes com todo o treino, destaca-se sobretudo o \textbf{SVM} neste ensaio.

\end{document}
